\documentclass[aspectratio=169,10pt]{beamer}
\usetheme{Madrid}
\usecolortheme{default}
\usepackage[utf8]{inputenc}
\usepackage[T1]{fontenc}
\usepackage{lmodern}
\usepackage{amsmath,amssymb}
\usepackage{enumitem}
\usepackage{xcolor}

\setbeamertemplate{navigation symbols}{}
\setbeamertemplate{footline}[frame number]
\setbeamertemplate{itemize items}[circle]
\setbeamerfont{frametitle}{size=\large,series=\bfseries}

\newcommand{\qitem}[4]{%
	\begin{frame}[t]{#1}
		\small
		\textbf{Question:} #2
		\vspace{4pt}
		\begin{enumerate}[label=\Alph*.,itemsep=1pt,topsep=2pt]
			#3
		\end{enumerate}
		\vspace{4pt}
		\textbf{Correct Answer:} #4
	\end{frame}
}

\newcommand{\expl}[1]{%
	\begin{frame}[t]{Explanation}
		\small #1
	\end{frame}
}

\title{CIA Part 1 -- 100 Hard Questions with Detailed Explanations}
\subtitle{Essentials of Internal Auditing}
\author{Study Deck}
\date{\today}

\begin{document}
	
	\begin{frame}
		\titlepage
	\end{frame}
	
	\begin{frame}{Format}
		\begin{itemize}
			\item 100 question slides (one question per slide).
			\item Each slide: question, options A--D, correct answer.
			\item Detailed explanation on the following slide (why correct, why others wrong).
			\item Domains: Foundations, Independence \& Objectivity, Proficiency \& Due Care,
			QAIP, Governance, Risk, Control, Fraud, Engagement Planning \& Reporting.
		\end{itemize}
	\end{frame}
	
	% ============ Q1 ============
	\begin{frame}[t]{Q1}
		\small
		\textbf{Question:} Which statement best reflects the IIA definition of internal auditing?
		\begin{enumerate}[label=\Alph*.,itemsep=1pt,topsep=2pt]
			\item A function that guarantees financial statements are free of misstatement.
			\item An independent, objective assurance and consulting activity designed to add value and improve operations.
			\item A compliance unit reporting to the external auditor.
			\item A management advisory service focused only on cost reduction.
		\end{enumerate}
		\vspace{4pt}
		\textbf{Correct Answer: B}
	\end{frame}
	
	\begin{frame}[t]{Q1 -- Explanation}
		\small
		\textbf{Correct: B.} The IIA definition states internal auditing is an independent, objective assurance and consulting activity designed to add value and improve an organization's operations. It helps achieve objectives through a systematic, disciplined approach to evaluating and improving risk management, control, and governance. \\[4pt]
		\textbf{A wrong:} Internal audit provides reasonable assurance, not a guarantee. Only external auditors opine on financial statements. \\[2pt]
		\textbf{C wrong:} Internal audit reports to the board/audit committee, not the external auditor, and is broader than compliance. \\[2pt]
		\textbf{D wrong:} Cost reduction is only one possible outcome; the scope is risk, control, and governance, not just cost.
	\end{frame}
	
	% ============ Q2 ============
	\begin{frame}[t]{Q2}
		\small
		\textbf{Question:} The primary purpose of the internal audit charter is to:
		\begin{enumerate}[label=\Alph*.,itemsep=1pt,topsep=2pt]
			\item Describe the annual audit plan.
			\item Define the purpose, authority, and responsibility of the internal audit activity.
			\item List the qualifications of each auditor.
			\item Establish the audit committee's meeting schedule.
		\end{enumerate}
		\vspace{4pt}
		\textbf{Correct Answer: B}
	\end{frame}
	
	\begin{frame}[t]{Q2 -- Explanation}
		\small
		\textbf{Correct: B.} Standard 1000 requires the purpose, authority, and responsibility of the internal audit activity to be formally defined in a charter, approved by the board. The charter establishes the activity's position and reporting lines. \\[4pt]
		\textbf{A wrong:} The annual audit plan is a separate document, risk-based and approved by the board. \\[2pt]
		\textbf{C wrong:} Auditor qualifications belong in HR records or the QAIP, not the charter. \\[2pt]
		\textbf{D wrong:} The audit committee's schedule is set by the committee, not the internal audit charter.
	\end{frame}
	
	% ============ Q3 ============
	\begin{frame}[t]{Q3}
		\small
		\textbf{Question:} Which reporting line best preserves internal audit independence?
		\begin{enumerate}[label=\Alph*.,itemsep=1pt,topsep=2pt]
			\item Directly to the controller.
			\item Functionally to the audit committee, administratively to the CEO.
			\item To the external auditor.
			\item To the CFO.
		\end{enumerate}
		\vspace{4pt}
		\textbf{Correct Answer: B}
	\end{frame}
	
	\begin{frame}[t]{Q3 -- Explanation}
		\small
		\textbf{Correct: B.} Functional reporting to the audit committee (charter, plan, reporting, performance) preserves independence. Administrative reporting to the CEO handles day-to-day HR, budget, and resources. This dual-line structure is best practice per IIA Standards. \\[4pt]
		\textbf{A wrong:} Reporting to the controller places internal audit under the very function it audits. \\[2pt]
		\textbf{C wrong:} Internal audit must not report to the external auditor; they are separate functions. \\[2pt]
		\textbf{D wrong:} Reporting to the CFO impairs independence over financial reporting audits.
	\end{frame}
	
	% ============ Q4 ============
	\begin{frame}[t]{Q4}
		\small
		\textbf{Question:} Which situation most clearly impairs an internal auditor's objectivity?
		\begin{enumerate}[label=\Alph*.,itemsep=1pt,topsep=2pt]
			\item Auditing a process the auditor designed two years ago.
			\item Auditing a department the auditor has never worked in.
			\item Using a standardized audit program.
			\item Reviewing a peer's working papers.
		\end{enumerate}
		\vspace{4pt}
		\textbf{Correct Answer: A}
	\end{frame}
	
	\begin{frame}[t]{Q4 -- Explanation}
		\small
		\textbf{Correct: A.} Objectivity is impaired when an auditor audits an activity they previously designed, implemented, or operated. Standard 1130 requires reassignment if objectivity is impaired. A two-year gap does not remove the impairment. \\[4pt]
		\textbf{B wrong:} No prior involvement = no impairment. \\[2pt]
		\textbf{C wrong:} Standardized programs improve consistency; they don't impair objectivity. \\[2pt]
		\textbf{D wrong:} Peer review is a normal QAIP activity and supports objectivity.
	\end{frame}
	
	% ============ Q5 ============
	\begin{frame}[t]{Q5}
		\small
		\textbf{Question:} Due professional care requires an internal auditor to:
		\begin{enumerate}[label=\Alph*.,itemsep=1pt,topsep=2pt]
			\item Detect all fraud.
			\item Exercise the care expected of a reasonably prudent auditor.
			\item Guarantee that controls are effective.
			\item Eliminate all risk.
		\end{enumerate}
		\vspace{4pt}
		\textbf{Correct Answer: B}
	\end{frame}
	
	\begin{frame}[t]{Q5 -- Explanation}
		\small
		\textbf{Correct: B.} Standard 1220: auditors must apply the care and skill expected of a reasonably prudent and competent auditor. It is not a guarantee -- it is reasonable care. \\[4pt]
		\textbf{A wrong:} No auditor can detect all fraud; due care reduces but doesn't eliminate risk. \\[2pt]
		\textbf{C wrong:} Internal audit provides reasonable assurance, not a guarantee. \\[2pt]
		\textbf{D wrong:} Risk cannot be eliminated; it can only be managed.
	\end{frame}
	
	% ============ Q6 ============
	\begin{frame}[t]{Q6}
		\small
		\textbf{Question:} Continuing professional development (CPD) is required for internal auditors primarily to:
		\begin{enumerate}[label=\Alph*.,itemsep=1pt,topsep=2pt]
			\item Satisfy HR paperwork.
			\item Maintain proficiency and competence.
			\item Increase billable hours.
			\item Reduce audit fees.
		\end{enumerate}
		\vspace{4pt}
		\textbf{Correct Answer: B}
	\end{frame}
	
	\begin{frame}[t]{Q6 -- Explanation}
		\small
		\textbf{Correct: B.} Standard 1230 requires internal auditors to enhance knowledge, skills, and competencies through CPD. This keeps the activity proficient. \\[4pt]
		\textbf{A wrong:} HR records are a byproduct, not the purpose. \\[2pt]
		\textbf{C wrong:} Billable hours are irrelevant to the Standards. \\[2pt]
		\textbf{D wrong:} Fee reduction is not the objective of CPD.
	\end{frame}
	
	% ============ Q7 ============
	\begin{frame}[t]{Q7}
		\small
		\textbf{Question:} The QAIP must include:
		\begin{enumerate}[label=\Alph*.,itemsep=1pt,topsep=2pt]
			\item Only internal assessments.
			\item Only external assessments.
			\item Both internal and external assessments.
			\item Neither, if the CAE is certified.
		\end{enumerate}
		\vspace{4pt}
		\textbf{Correct Answer: C}
	\end{frame}
	
	\begin{frame}[t]{Q7 -- Explanation}
		\small
		\textbf{Correct: C.} Standard 1300 requires a QAIP covering all aspects of the internal audit activity, including both ongoing/periodic internal assessments and an external assessment at least every five years. \\[4pt]
		\textbf{A wrong:} Internal-only misses the required external assessment. \\[2pt]
		\textbf{B wrong:} External-only misses ongoing internal monitoring. \\[2pt]
		\textbf{D wrong:} Certification of the CAE does not waive QAIP requirements.
	\end{frame}
	
	% ============ Q8 ============
	\begin{frame}[t]{Q8}
		\small
		\textbf{Question:} An external quality assessment must be conducted at least every:
		\begin{enumerate}[label=\Alph*.,itemsep=1pt,topsep=2pt]
			\item Two years.
			\item Three years.
			\item Five years.
			\item Seven years.
		\end{enumerate}
		\vspace{4pt}
		\textbf{Correct Answer: C}
	\end{frame}
	
	\begin{frame}[t]{Q8 -- Explanation}
		\small
		\textbf{Correct: C.} Standard 1312 requires an external assessment at least once every five years by a qualified, independent reviewer or team from outside the organization. \\[4pt]
		\textbf{A wrong:} Two years is too frequent and not required. \\[2pt]
		\textbf{B wrong:} Three years is not the Standard. \\[2pt]
		\textbf{D wrong:} Seven years exceeds the maximum interval.
	\end{frame}
	
	% ============ Q9 ============
	\begin{frame}[t]{Q9}
		\small
		\textbf{Question:} Which is NOT one of the four principles of the IIA Code of Ethics?
		\begin{enumerate}[label=\Alph*.,itemsep=1pt,topsep=2pt]
			\item Integrity.
			\item Objectivity.
			\item Profitability.
			\item Confidentiality.
		\end{enumerate}
		\vspace{4pt}
		\textbf{Correct Answer: C}
	\end{frame}
	
	\begin{frame}[t]{Q9 -- Explanation}
		\small
		\textbf{Correct: C.} The four principles are Integrity, Objectivity, Confidentiality, and Competency. Profitability is not a principle. \\[4pt]
		\textbf{A wrong:} Integrity is required -- honesty and candor. \\[2pt]
		\textbf{B wrong:} Objectivity is required -- no conflicts, no bias. \\[2pt]
		\textbf{D wrong:} Confidentiality is required -- protect information.
	\end{frame}
	
	% ============ Q10 ============
	\begin{frame}[t]{Q10}
		\small
		\textbf{Question:} Governance is best defined as:
		\begin{enumerate}[label=\Alph*.,itemsep=1pt,topsep=2pt]
			\item Preparing financial statements.
			\item The combination of processes and structures used to direct and control the organization.
			\item The internal audit plan.
			\item The external audit opinion.
		\end{enumerate}
		\vspace{4pt}
		\textbf{Correct Answer: B}
	\end{frame}
	
	\begin{frame}[t]{Q10 -- Explanation}
		\small
		\textbf{Correct: B.} Governance is the set of processes and structures (board, policies, accountability) that direct and control the organization and align it with stakeholder interests. \\[4pt]
		\textbf{A wrong:} Financial statement preparation is a management function under governance. \\[2pt]
		\textbf{C wrong:} The audit plan is a tool of internal audit, not governance itself. \\[2pt]
		\textbf{D wrong:} The external audit opinion is an output, not governance.
	\end{frame}
	
	% ============ Q11 ============
	\begin{frame}[t]{Q11}
		\small
		\textbf{Question:} Inherent risk is:
		\begin{enumerate}[label=\Alph*.,itemsep=1pt,topsep=2pt]
			\item Risk after controls.
			\item Risk before considering controls.
			\item Risk management ignores.
			\item Risk that is always material.
		\end{enumerate}
		\vspace{4pt}
		\textbf{Correct Answer: B}
	\end{frame}
	
	\begin{frame}[t]{Q11 -- Explanation}
		\small
		\textbf{Correct: B.} Inherent risk is the risk of an event occurring before any controls are applied. It is a component of the audit risk model along with control and detection risk. \\[4pt]
		\textbf{A wrong:} That describes residual risk. \\[2pt]
		\textbf{C wrong:} No risk should be ignored; inherent risk is assessed, not dismissed. \\[2pt]
		\textbf{D wrong:} Inherent risk can be high or low; materiality is separate.
	\end{frame}
	
	% ============ Q12 ============
	\begin{frame}[t]{Q12}
		\small
		\textbf{Question:} Residual risk is:
		\begin{enumerate}[label=\Alph*.,itemsep=1pt,topsep=2pt]
			\item Risk before controls.
			\item Risk remaining after controls.
			\item Risk of fraud only.
			\item Risk of auditor error.
		\end{enumerate}
		\vspace{4pt}
		\textbf{Correct Answer: B}
	\end{frame}
	
	\begin{frame}[t]{Q12 -- Explanation}
		\small
		\textbf{Correct: B.} Residual risk = inherent risk minus the effect of controls. It is the risk that remains after management's response. \\[4pt]
		\textbf{A wrong:} That is inherent risk. \\[2pt]
		\textbf{C wrong:} Fraud is one type of risk, not residual risk. \\[2pt]
		\textbf{D wrong:} Auditor error is covered by detection risk, a different concept.
	\end{frame}
	
	% ============ Q13 ============
	\begin{frame}[t]{Q13}
		\small
		\textbf{Question:} Which is a preventive control?
		\begin{enumerate}[label=\Alph*.,itemsep=1pt,topsep=2pt]
			\item Bank reconciliation.
			\item Segregation of duties.
			\item Variance analysis.
			\item Exception report.
		\end{enumerate}
		\vspace{4pt}
		\textbf{Correct Answer: B}
	\end{frame}
	
	\begin{frame}[t]{Q13 -- Explanation}
		\small
		\textbf{Correct: B.} Segregation of duties prevents fraud/error by stopping one person from controlling all aspects of a transaction. It acts before the event. \\[4pt]
		\textbf{A wrong:} Reconciliation detects errors after the fact. \\[2pt]
		\textbf{C wrong:} Variance analysis is detective. \\[2pt]
		\textbf{D wrong:} Exception reports are detective.
	\end{frame}
	
	% ============ Q14 ============
	\begin{frame}[t]{Q14}
		\small
		\textbf{Question:} Which is a detective control?
		\begin{enumerate}[label=\Alph*.,itemsep=1pt,topsep=2pt]
			\item Authorization.
			\item Physical safeguards.
			\item Bank reconciliation.
			\item Password policy.
		\end{enumerate}
		\vspace{4pt}
		\textbf{Correct Answer: C}
	\end{frame}
	
	\begin{frame}[t]{Q14 -- Explanation}
		\small
		\textbf{Correct: C.} Reconciliation detects discrepancies after they occur, enabling correction. \\[4pt]
		\textbf{A wrong:} Authorization is preventive. \\[2pt]
		\textbf{B wrong:} Physical safeguards are preventive. \\[2pt]
		\textbf{D wrong:} Password policy is preventive.
	\end{frame}
	
	% ============ Q15 ============
	\begin{frame}[t]{Q15}
		\small
		\textbf{Question:} The fraud triangle consists of:
		\begin{enumerate}[label=\Alph*.,itemsep=1pt,topsep=2pt]
			\item Pressure, opportunity, rationalization.
			\item Risk, control, detection.
			\item Plan, do, check.
			\item Cash, inventory, receivables.
		\end{enumerate}
		\vspace{4pt}
		\textbf{Correct Answer: A}
	\end{frame}
	
	\begin{frame}[t]{Q15 -- Explanation}
		\small
		\textbf{Correct: A.} Cressey's fraud triangle: pressure (incentive), opportunity, and rationalization. All three usually must be present for fraud. \\[4pt]
		\textbf{B wrong:} That's the audit risk model. \\[2pt]
		\textbf{C wrong:} That's the Deming/PDCA cycle. \\[2pt]
		\textbf{D wrong:} Those are balance sheet accounts, not a fraud model.
	\end{frame}
	
	% ============ Q16 ============
	\begin{frame}[t]{Q16}
		\small
		\textbf{Question:} When fraud is suspected, the internal auditor should:
		\begin{enumerate}[label=\Alph*.,itemsep=1pt,topsep=2pt]
			\item Confront the suspect immediately.
			\item Notify appropriate authorities within the organization and recommend investigation.
			\item Ignore it until year-end.
			\item Destroy workpapers.
		\end{enumerate}
		\vspace{4pt}
		\textbf{Correct Answer: B}
	\end{frame}
	
	\begin{frame}[t]{Q16 -- Explanation}
		\small
		\textbf{Correct: B.} Standard 1210.A2: auditors must have sufficient knowledge to identify fraud indicators and report them to appropriate authorities; investigation is management/legal's job. \\[4pt]
		\textbf{A wrong:} Confronting the suspect can tip them off and destroy evidence. \\[2pt]
		\textbf{C wrong:} Delay allows fraud to continue and evidence to vanish. \\[2pt]
		\textbf{D wrong:} Destroying workpapers is unethical and potentially illegal.
	\end{frame}
	
	% ============ Q17 ============
	\begin{frame}[t]{Q17}
		\small
		\textbf{Question:} An audit finding should include:
		\begin{enumerate}[label=\Alph*.,itemsep=1pt,topsep=2pt]
			\item Criteria, condition, cause, effect, recommendation.
			\item Only the recommendation.
			\item Only the condition.
			\item Management's response only.
		\end{enumerate}
		\vspace{4pt}
		\textbf{Correct Answer: A}
	\end{frame}
	
	\begin{frame}[t]{Q17 -- Explanation}
		\small
		\textbf{Correct: A.} The five-element finding model: Criteria (what should be), Condition (what is), Cause (why), Effect (impact), and Recommendation (fix). \\[4pt]
		\textbf{B wrong:} Recommendation alone lacks context. \\[2pt]
		\textbf{C wrong:} Condition alone lacks criteria and impact. \\[2pt]
		\textbf{D wrong:} Management's response is part of the report but not the finding itself.
	\end{frame}
	
	% ============ Q18 ============
	\begin{frame}[t]{Q18}
		\small
		\textbf{Question:} Audit risk is the risk that:
		\begin{enumerate}[label=\Alph*.,itemsep=1pt,topsep=2pt]
			\item The client goes bankrupt.
			\item The auditor reaches an incorrect conclusion.
			\item Controls fail.
			\item Fraud occurs.
		\end{enumerate}
		\vspace{4pt}
		\textbf{Correct Answer: B}
	\end{frame}
	
	\begin{frame}[t]{Q18 -- Explanation}
		\small
		\textbf{Correct: B.} Audit risk is the risk the auditor issues an inappropriate conclusion. It is modeled as AR = IR x CR x DR. \\[4pt]
		\textbf{A wrong:} Bankruptcy is business risk. \\[2pt]
		\textbf{C wrong:} Control failure contributes to control risk, not audit risk definition. \\[2pt]
		\textbf{D wrong:} Fraud is a source of risk, not the definition of audit risk.
	\end{frame}
	
	% ============ Q19 ============
	\begin{frame}[t]{Q19}
		\small
		\textbf{Question:} Which is the strongest form of audit evidence?
		\begin{enumerate}[label=\Alph*.,itemsep=1pt,topsep=2pt]
			\item Inquiry.
			\item Observation.
			\item Physical examination by the auditor.
			\item Analytical procedures.
		\end{enumerate}
		\vspace{4pt}
		\textbf{Correct Answer: C}
	\end{frame}
	
	\begin{frame}[t]{Q19 -- Explanation}
		\small
		\textbf{Correct: C.} Evidence directly obtained by the auditor through physical examination is most reliable because the auditor controls its collection. \\[4pt]
		\textbf{A wrong:} Inquiry is the weakest -- it relies on others' representations. \\[2pt]
		\textbf{B wrong:} Observation is good but limited to the moment of observation. \\[2pt]
		\textbf{D wrong:} Analytical procedures are indirect and depend on data quality.
	\end{frame}
	
	% ============ Q20 ============
	\begin{frame}[t]{Q20}
		\small
		\textbf{Question:} Sufficiency of evidence relates to:
		\begin{enumerate}[label=\Alph*.,itemsep=1pt,topsep=2pt]
			\item Quality.
			\item Quantity.
			\item Timeliness.
			\item Relevance.
		\end{enumerate}
		\vspace{4pt}
		\textbf{Correct Answer: B}
	\end{frame}
	
	\begin{frame}[t]{Q20 -- Explanation}
		\small
		\textbf{Correct: B.} Sufficiency = quantity of evidence. Appropriateness = relevance + reliability (quality). \\[4pt]
		\textbf{A wrong:} Quality is appropriateness. \\[2pt]
		\textbf{C wrong:} Timeliness is a dimension of reliability, not sufficiency. \\[2pt]
		\textbf{D wrong:} Relevance is part of appropriateness.
	\end{frame}
	
	% ============ Q21 ============
	\begin{frame}[t]{Q21}
		\small
		\textbf{Question:} Which sampling method gives every item an equal chance of selection?
		\begin{enumerate}[label=\Alph*.,itemsep=1pt,topsep=2pt]
			\item Judgmental sampling.
			\item Haphazard sampling.
			\item Random sampling.
			\item Block sampling.
		\end{enumerate}
		\vspace{4pt}
		\textbf{Correct Answer: C}
	\end{frame}
	
	\begin{frame}[t]{Q21 -- Explanation}
		\small
		\textbf{Correct: C.} Random sampling uses a random number generator or table so every item has a known, equal chance. It supports statistical inference. \\[4pt]
		\textbf{A wrong:} Judgmental relies on auditor bias. \\[2pt]
		\textbf{B wrong:} Haphazard is not statistically valid. \\[2pt]
		\textbf{D wrong:} Block sampling selects contiguous items and is biased.
	\end{frame}
	
	% ============ Q22 ============
	\begin{frame}[t]{Q22}
		\small
		\textbf{Question:} Attribute sampling is used to:
		\begin{enumerate}[label=\Alph*.,itemsep=1pt,topsep=2pt]
			\item Estimate a monetary amount.
			\item Estimate a rate of deviation from a control.
			\item Project fraud losses.
			\item Set materiality.
		\end{enumerate}
		\vspace{4pt}
		\textbf{Correct Answer: B}
	\end{frame}
	
	\begin{frame}[t]{Q22 -- Explanation}
		\small
		\textbf{Correct: B.} Attribute sampling tests whether a characteristic (attribute) exists -- e.g., was the invoice approved? It estimates deviation rate for control testing. \\[4pt]
		\textbf{A wrong:} Monetary unit sampling estimates amounts. \\[2pt]
		\textbf{C wrong:} Fraud loss projection uses different techniques. \\[2pt]
		\textbf{D wrong:} Materiality is a judgment, not a sampling output.
	\end{frame}
	
	% ============ Q23 ============
	\begin{frame}[t]{Q23}
		\small
		\textbf{Question:} The first step in an audit engagement is to:
		\begin{enumerate}[label=\Alph*.,itemsep=1pt,topsep=2pt]
			\item Perform fieldwork.
			\item Issue the report.
			\item Define objectives and scope.
			\item Conduct the exit meeting.
		\end{enumerate}
		\vspace{4pt}
		\textbf{Correct Answer: C}
	\end{frame}
	
	\begin{frame}[t]{Q23 -- Explanation}
		\small
		\textbf{Correct: C.} Planning begins with defining engagement objectives, scope, and criteria. Everything else flows from that. \\[4pt]
		\textbf{A wrong:} Fieldwork comes after planning. \\[2pt]
		\textbf{B wrong:} The report is the last step. \\[2pt]
		\textbf{D wrong:} The exit meeting occurs near the end.
	\end{frame}
	
	% ============ Q24 ============
	\begin{frame}[t]{Q24}
		\small
		\textbf{Question:} A consulting engagement differs from an assurance engagement because consulting:
		\begin{enumerate}[label=\Alph*.,itemsep=1pt,topsep=2pt]
			\item Always produces a formal opinion.
			\item Is advisory in nature and typically requested by the client.
			\item Cannot be performed by internal audit.
			\item Requires no documentation.
		\end{enumerate}
		\vspace{4pt}
		\textbf{Correct Answer: B}
	\end{frame}
	
	\begin{frame}[t]{Q24 -- Explanation}
		\small
		\textbf{Correct: B.} Consulting is advisory and usually requested by management. Assurance is an objective assessment where internal audit determines scope. \\[4pt]
		\textbf{A wrong:} Assurance, not consulting, produces conclusions/opinions. \\[2pt]
		\textbf{C wrong:} Internal audit may perform consulting per its charter. \\[2pt]
		\textbf{D wrong:} All engagements must be documented.
	\end{frame}
	
	% ============ Q25 ============
	\begin{frame}[t]{Q25}
		\small
		\textbf{Question:} Before accepting a consulting engagement, the CAE should first consider:
		\begin{enumerate}[label=\Alph*.,itemsep=1pt,topsep=2pt]
			\item Whether it impairs objectivity or independence.
			\item Only the fee.
			\item Only the client's preference.
			\item Only the timeline.
		\end{enumerate}
		\vspace{4pt}
		\textbf{Correct Answer: A}
	\end{frame}
	
	\begin{frame}[t]{Q25 -- Explanation}
		\small
		\textbf{Correct: A.} Standard 1130.C1: the CAE may accept consulting engagements, but must decline if objectivity or independence would be impaired. \\[4pt]
		\textbf{B wrong:} Fees are not the primary consideration. \\[2pt]
		\textbf{C wrong:} Client preference does not override the Standards. \\[2pt]
		\textbf{D wrong:} Timeline is secondary to independence.
	\end{frame}
	
	% ============ Q26 ============
	\begin{frame}[t]{Q26}
		\small
		\textbf{Question:} Which element is NOT part of the audit risk model?
		\begin{enumerate}[label=\Alph*.,itemsep=1pt,topsep=2pt]
			\item Inherent risk.
			\item Control risk.
			\item Detection risk.
			\item Market risk.
		\end{enumerate}
		\vspace{4pt}
		\textbf{Correct Answer: D}
	\end{frame}
	
	\begin{frame}[t]{Q26 -- Explanation}
		\small
		\textbf{Correct: D.} The audit risk model is AR = IR x CR x DR. Market risk is a business/economic concept, not part of the audit risk model. \\[4pt]
		\textbf{A wrong:} IR is part of the model. \\[2pt]
		\textbf{B wrong:} CR is part of the model. \\[2pt]
		\textbf{C wrong:} DR is part of the model.
	\end{frame}
	
	% ============ Q27 ============
	\begin{frame}[t]{Q27}
		\small
		\textbf{Question:} Materiality in an audit is:
		\begin{enumerate}[label=\Alph*.,itemsep=1pt,topsep=2pt]
			\item A fixed percentage of revenue.
			\item A matter of professional judgment about what influences decisions.
			\item Always \$1 million.
			\item Set only by the external auditor.
		\end{enumerate}
		\vspace{4pt}
		\textbf{Correct Answer: B}
	\end{frame}
	
	\begin{frame}[t]{Q27 -- Explanation}
		\small
		\textbf{Correct: B.} Materiality is professional judgment: would the omission/misstatement influence a reasonable user's decisions? It varies by entity and context. \\[4pt]
		\textbf{A wrong:} No fixed percentage is mandated. \\[2pt]
		\textbf{C wrong:} No universal dollar threshold exists. \\[2pt]
		\textbf{D wrong:} Internal auditors also assess materiality for their engagements.
	\end{frame}
	
	% ============ Q28 ============
	\begin{frame}[t]{Q28}
		\small
		\textbf{Question:} A walkthrough is used primarily to:
		\begin{enumerate}[label=\Alph*.,itemsep=1pt,topsep=2pt]
			\item Confirm account balances.
			\item Understand the flow of a transaction through the system.
			\item Test IT access controls.
			\item Calculate materiality.
		\end{enumerate}
		\vspace{4pt}
		\textbf{Correct Answer: B}
	\end{frame}
	
	\begin{frame}[t]{Q28 -- Explanation}
		\small
		\textbf{Correct: B.} A walkthrough traces one or a few transactions from start to finish to understand process flow, controls, and gaps. \\[4pt]
		\textbf{A wrong:} Confirmations verify balances, not process flow. \\[2pt]
		\textbf{C wrong:} IT access testing is a separate procedure. \\[2pt]
		\textbf{D wrong:} Materiality is a judgment, not a walkthrough output.
	\end{frame}
	
	% ============ Q29 ============
	\begin{frame}[t]{Q29}
		\small
		\textbf{Question:} The best evidence of the existence of inventory is:
		\begin{enumerate}[label=\Alph*.,itemsep=1pt,topsep=2pt]
			\item Management's representation.
			\item Physical observation/count by the auditor.
			\item The prior year's workpapers.
			\item The purchasing manager's email.
		\end{enumerate}
		\vspace{4pt}
		\textbf{Correct Answer: B}
	\end{frame}
	
	\begin{frame}[t]{Q29 -- Explanation}
		\small
		\textbf{Correct: B.} Direct physical observation/count by the auditor provides the strongest evidence of existence. \\[4pt]
		\textbf{A wrong:} Management's representation is weak and self-serving. \\[2pt]
		\textbf{C wrong:} Prior-year workpapers do not prove current existence. \\[2pt]
		\textbf{D wrong:} An email is documentary evidence of low reliability.
	\end{frame}
	
	% ============ Q30 ============
	\begin{frame}[t]{Q30}
		\small
		\textbf{Question:} Which is the best way to test the operating effectiveness of a control?
		\begin{enumerate}[label=\Alph*.,itemsep=1pt,topsep=2pt]
			\item Inquiry.
			\item Observation.
			\item Reperformance.
			\item Reading the policy.
		\end{enumerate}
		\vspace{4pt}
		\textbf{Correct Answer: C}
	\end{frame}
	
	\begin{frame}[t]{Q30 -- Explanation}
		\small
		\textbf{Correct: C.} Reperformance -- the auditor independently executes the control -- provides direct, high-quality evidence of operating effectiveness. \\[4pt]
		\textbf{A wrong:} Inquiry alone is insufficient. \\[2pt]
		\textbf{B wrong:} Observation is limited to the observation period. \\[2pt]
		\textbf{D wrong:} Reading policy tests design, not operation.
	\end{frame}
	
	% ============ Q31 ============
	\begin{frame}[t]{Q31}
		\small
		\textbf{Question:} The CAE's functional reporting line to the board ensures:
		\begin{enumerate}[label=\Alph*.,itemsep=1pt,topsep=2pt]
			\item Budget approval.
			\item Independence and authority.
			\item Daily oversight.
			\item HR administration.
		\end{enumerate}
		\vspace{4pt}
		\textbf{Correct Answer: B}
	\end{frame}
	
	\begin{frame}[t]{Q31 -- Explanation}
		\small
		\textbf{Correct: B.} Functional reporting to the board protects independence and gives internal audit the authority to act. \\[4pt]
		\textbf{A wrong:} Budget approval is administrative. \\[2pt]
		\textbf{C wrong:} Daily oversight is management's job. \\[2pt]
		\textbf{D wrong:} HR is administrative, handled by the CEO line.
	\end{frame}
	
	% ============ Q32 ============
	\begin{frame}[t]{Q32}
		\small
		\textbf{Question:} Which best describes "tone at the top"?
		\begin{enumerate}[label=\Alph*.,itemsep=1pt,topsep=2pt]
			\item The CEO's salary.
			\item Management's commitment to ethics and control.
			\item The marketing slogan.
			\item The IT password policy.
		\end{enumerate}
		\vspace{4pt}
		\textbf{Correct Answer: B}
	\end{frame}
	
	\begin{frame}[t]{Q32 -- Explanation}
		\small
		\textbf{Correct: B.} Tone at the top is the control environment element reflecting leadership's ethics, integrity, and commitment to control. \\[4pt]
		\textbf{A wrong:} Compensation is unrelated. \\[2pt]
		\textbf{C wrong:} Marketing slogans are not control environment. \\[2pt]
		\textbf{D wrong:} Password policy is one IT control, not tone at the top.
	\end{frame}
	
	% ============ Q33 ============
	\begin{frame}[t]{Q33}
		\small
		\textbf{Question:} A control deficiency exists when:
		\begin{enumerate}[label=\Alph*.,itemsep=1pt,topsep=2pt]
			\item Controls are documented.
			\item Design or operation does not prevent/detect misstatements timely.
			\item Management disagrees.
			\item Risk is zero.
		\end{enumerate}
		\vspace{4pt}
		\textbf{Correct Answer: B}
	\end{frame}
	
	\begin{frame}[t]{Q33 -- Explanation}
		\small
		\textbf{Correct: B.} A deficiency exists when a control is missing, poorly designed, or not operating so that misstatements are not prevented/detected timely. \\[4pt]
		\textbf{A wrong:} Documentation is a strength, not a deficiency. \\[2pt]
		\textbf{C wrong:} Management disagreement doesn't create a deficiency. \\[2pt]
		\textbf{D wrong:} Zero risk is unrealistic and not a deficiency test.
	\end{frame}
	
	% ============ Q34 ============
	\begin{frame}[t]{Q34}
		\small
		\textbf{Question:} A material weakness is:
		\begin{enumerate}[label=\Alph*.,itemsep=1pt,topsep=2pt]
			\item A minor deficiency.
			\item A deficiency creating a reasonable possibility of material misstatement.
			\item A documentation gap.
			\item A training issue.
		\end{enumerate}
		\vspace{4pt}
		\textbf{Correct Answer: B}
	\end{frame}
	
	\begin{frame}[t]{Q34 -- Explanation}
		\small
		\textbf{Correct: B.} A material weakness is a deficiency (or combination) that creates a reasonable possibility that a material misstatement will not be prevented or detected timely. \\[4pt]
		\textbf{A wrong:} Minor deficiencies are less severe. \\[2pt]
		\textbf{C wrong:} Documentation gaps may be deficiencies but not automatically material. \\[2pt]
		\textbf{D wrong:} Training issues may contribute but are not the definition.
	\end{frame}
	
	% ============ Q35 ============
	\begin{frame}[t]{Q35}
		\small
		\textbf{Question:} Risk appetite is:
		\begin{enumerate}[label=\Alph*.,itemsep=1pt,topsep=2pt]
			\item The amount of risk an organization is willing to accept.
			\item The amount of risk the auditor accepts.
			\item Always zero.
			\item Set by regulators.
		\end{enumerate}
		\vspace{4pt}
		\textbf{Correct Answer: A}
	\end{frame}
	
	\begin{frame}[t]{Q35 -- Explanation}
		\small
		\textbf{Correct: A.} Risk appetite is the amount and type of risk an organization is willing to pursue or retain in pursuit of objectives. \\[4pt]
		\textbf{B wrong:} Auditor risk acceptance is audit risk, a different concept. \\[2pt]
		\textbf{C wrong:} Zero appetite is unrealistic. \\[2pt]
		\textbf{D wrong:} Regulators set minimums, not the entity's appetite.
	\end{frame}
	
	% ============ Q36 ============
	\begin{frame}[t]{Q36}
		\small
		\textbf{Question:} Which is an example of a consulting engagement?
		\begin{enumerate}[label=\Alph*.,itemsep=1pt,topsep=2pt]
			\item Financial statement audit.
			\item Compliance audit.
			\item Advisory on a new ERP control design.
			\item Operational audit.
		\end{enumerate}
		\vspace{4pt}
		\textbf{Correct Answer: C}
	\end{frame}
	
	\begin{frame}[t]{Q36 -- Explanation}
		\small
		\textbf{Correct: C.} Advising on control design for a new ERP is advisory (consulting). Assurance engagements involve objective assessment of existing conditions. \\[4pt]
		\textbf{A wrong:} Financial statement audit is external assurance. \\[2pt]
		\textbf{B wrong:} Compliance audit is assurance. \\[2pt]
		\textbf{D wrong:} Operational audit is assurance.
	\end{frame}
	
	% ============ Q37 ============
	\begin{frame}[t]{Q37}
		\small
		\textbf{Question:} An exit conference is used to:
		\begin{enumerate}[label=\Alph*.,itemsep=1pt,topsep=2pt]
			\item Begin fieldwork.
			\item Discuss findings and management responses before the final report.
			\item Replace the audit report.
			\item Set the audit fee.
		\end{enumerate}
		\vspace{4pt}
		\textbf{Correct Answer: B}
	\end{frame}
	
	\begin{frame}[t]{Q37 -- Explanation}
		\small
		\textbf{Correct: B.} The exit conference validates findings and obtains management's response before issuing the final report, reducing misunderstandings. \\[4pt]
		\textbf{A wrong:} Fieldwork starts before the exit conference. \\[2pt]
		\textbf{C wrong:} It supplements, not replaces, the report. \\[2pt]
		\textbf{D wrong:} Fees are unrelated to the exit conference.
	\end{frame}
	
	% ============ Q38 ============
	\begin{frame}[t]{Q38}
		\small
		\textbf{Question:} The final engagement communication should include:
		\begin{enumerate}[label=\Alph*.,itemsep=1pt,topsep=2pt]
			\item Only negative findings.
			\item Objectives, scope, findings, conclusions, recommendations.
			\item Only management's response.
			\item Only the audit plan.
		\end{enumerate}
		\vspace{4pt}
		\textbf{Correct Answer: B}
	\end{frame}
	
	\begin{frame}[t]{Q38 -- Explanation}
		\small
		\textbf{Correct: B.} Standard 2400 requires the communication to include objectives, scope, findings, conclusions, and recommendations (where applicable). \\[4pt]
		\textbf{A wrong:} Positive observations may also be included. \\[2pt]
		\textbf{C wrong:} Management's response is one part, not the whole. \\[2pt]
		\textbf{D wrong:} The audit plan is not the communication.
	\end{frame}
	
	% ============ Q39 ============
	\begin{frame}[t]{Q39}
		\small
		\textbf{Question:} Follow-up on audit recommendations is:
		\begin{enumerate}[label=\Alph*.,itemsep=1pt,topsep=2pt]
			\item Optional.
			\item Required to determine whether corrective action was taken.
			\item The external auditor's job.
			\item Not part of internal audit.
		\end{enumerate}
		\vspace{4pt}
		\textbf{Correct Answer: B}
	\end{frame}
	
	\begin{frame}[t]{Q39 -- Explanation}
		\small
		\textbf{Correct: B.} Standard 2500 requires the CAE to establish a follow-up process to monitor and ensure management actions are effectively implemented. \\[4pt]
		\textbf{A wrong:} Follow-up is required, not optional. \\[2pt]
		\textbf{C wrong:} It is internal audit's responsibility. \\[2pt]
		\textbf{D wrong:} Follow-up is a core internal audit activity.
	\end{frame}
	
	% ============ Q40 ============
	\begin{frame}[t]{Q40}
		\small
		\textbf{Question:} The audit committee's primary responsibility is to:
		\begin{enumerate}[label=\Alph*.,itemsep=1pt,topsep=2pt]
			\item Manage the company.
			\item Oversee financial reporting, audit, and internal control.
			\item Set product strategy.
			\item Hire all staff.
		\end{enumerate}
		\vspace{4pt}
		\textbf{Correct Answer: B}
	\end{frame}
	
	\begin{frame}[t]{Q40 -- Explanation}
		\small
		\textbf{Correct: B.} The audit committee oversees financial reporting, external and internal audit, and internal control on behalf of the board. \\[4pt]
		\textbf{A wrong:} Management runs the company. \\[2pt]
		\textbf{C wrong:} Strategy is the full board/management's role. \\[2pt]
		\textbf{D wrong:} Hiring is management/HR.
	\end{frame}
	
	% ============ Q41 ============
	\begin{frame}[t]{Q41}
		\small
		\textbf{Question:} Which is the best evidence of a contract's terms?
		\begin{enumerate}[label=\Alph*.,itemsep=1pt,topsep=2pt]
			\item Oral confirmation.
			\item The signed contract itself.
			\item A legal department summary.
			\item An email referencing it.
		\end{enumerate}
		\vspace{4pt}
		\textbf{Correct Answer: B}
	\end{frame}
	
	\begin{frame}[t]{Q41 -- Explanation}
		\small
		\textbf{Correct: B.} The signed original contract is the primary, most reliable evidence of its terms. \\[4pt]
		\textbf{A wrong:} Oral evidence is the weakest. \\[2pt]
		\textbf{C wrong:} A summary is derivative. \\[2pt]
		\textbf{D wrong:} An email may not reflect final terms.
	\end{frame}
	
	% ============ Q42 ============
	\begin{frame}[t]{Q42}
		\small
		\textbf{Question:} Which would MOST likely impair auditor objectivity?
		\begin{enumerate}[label=\Alph*.,itemsep=1pt,topsep=2pt]
			\item Rotating audit assignments.
			\item Accepting a significant gift from an audit client.
			\item Using standardized work programs.
			\item Reviewing prior-year workpapers.
		\end{enumerate}
		\vspace{4pt}
		\textbf{Correct Answer: B}
	\end{frame}
	
	\begin{frame}[t]{Q42 -- Explanation}
		\small
		\textbf{Correct: B.} Accepting significant gifts creates a conflict of interest and impairs objectivity under the Code of Ethics. \\[4pt]
		\textbf{A wrong:} Rotation enhances objectivity. \\[2pt]
		\textbf{C wrong:} Standardized programs improve consistency. \\[2pt]
		\textbf{D wrong:} Prior workpapers provide context, not bias.
	\end{frame}
	
	% ============ Q43 ============
	\begin{frame}[t]{Q43}
		\small
		\textbf{Question:} The audit committee's role in internal audit includes:
		\begin{enumerate}[label=\Alph*.,itemsep=1pt,topsep=2pt]
			\item Approving the audit charter and plan.
			\item Performing fieldwork.
			\item Writing workpapers.
			\item Issuing audit reports.
		\end{enumerate}
		\vspace{4pt}
		\textbf{Correct Answer: A}
	\end{frame}
	
	\begin{frame}[t]{Q43 -- Explanation}
		\small
		\textbf{Correct: A.} The audit committee approves the internal audit charter, plan, budget, and appointment/removal of the CAE. \\[4pt]
		\textbf{B wrong:} Fieldwork is done by internal auditors. \\[2pt]
		\textbf{C wrong:} Workpapers are internal audit's. \\[2pt]
		\textbf{D wrong:} Reports are issued by internal audit.
	\end{frame}
	
	% ============ Q44 ============
	\begin{frame}[t]{Q44}
		\small
		\textbf{Question:} Which is NOT a component of the COSO Internal Control framework?
		\begin{enumerate}[label=\Alph*.,itemsep=1pt,topsep=2pt]
			\item Control environment.
			\item Risk assessment.
			\item Marketing strategy.
			\item Monitoring.
		\end{enumerate}
		\vspace{4pt}
		\textbf{Correct Answer: C}
	\end{frame}
	
	\begin{frame}[t]{Q44 -- Explanation}
		\small
		\textbf{Correct: C.} COSO's five components: Control Environment, Risk Assessment, Control Activities, Information \& Communication, Monitoring. Marketing strategy is not one. \\[4pt]
		\textbf{A wrong:} Control environment is COSO. \\[2pt]
		\textbf{B wrong:} Risk assessment is COSO. \\[2pt]
		\textbf{D wrong:} Monitoring is COSO.
	\end{frame}
	
	% ============ Q45 ============
	\begin{frame}[t]{Q45}
		\small
		\textbf{Question:} Which is the best description of fraud risk assessment?
		\begin{enumerate}[label=\Alph*.,itemsep=1pt,topsep=2pt]
			\item A one-time exercise.
			\item A continuous process integrated into audit planning.
			\item Only the external auditor's responsibility.
			\item Only relevant for public companies.
		\end{enumerate}
		\vspace{4pt}
		\textbf{Correct Answer: B}
	\end{frame}
	
	\begin{frame}[t]{Q45 -- Explanation}
		\small
		\textbf{Correct: B.} Fraud risk assessment is continuous, integrated into planning, and considers incentives, opportunities, and rationalization. \\[4pt]
		\textbf{A wrong:} One-time assessments miss emerging risks. \\[2pt]
		\textbf{C wrong:} Internal audit also assesses fraud risk. \\[2pt]
		\textbf{D wrong:} All organizations face fraud risk.
	\end{frame}
	
	% ============ Q46 ============
	\begin{frame}[t]{Q46}
		\small
		\textbf{Question:} A significant control deficiency is discovered. The CAE should FIRST:
		\begin{enumerate}[label=\Alph*.,itemsep=1pt,topsep=2pt]
			\item Report to the audit committee.
			\item Discuss with management to validate the finding.
			\item Notify external auditors.
			\item Issue a qualified opinion.
		\end{enumerate}
		\vspace{4pt}
		\textbf{Correct Answer: B}
	\end{frame}
	
	\begin{frame}[t]{Q46 -- Explanation}
		\small
		\textbf{Correct: B.} First validate the finding with management to confirm accuracy and get their response. Then escalate as needed. \\[4pt]
		\textbf{A wrong:} Reporting to the audit committee follows validation and management response. \\[2pt]
		\textbf{C wrong:} External auditors are notified later, if at all. \\[2pt]
		\textbf{D wrong:} Internal audit does not issue qualified opinions.
	\end{frame}
	
	% ============ Q47 ============
	\begin{frame}[t]{Q47}
		\small
		\textbf{Question:} An internal auditor lacks the expertise for an engagement. The BEST action is to:
		\begin{enumerate}[label=\Alph*.,itemsep=1pt,topsep=2pt]
			\item Proceed anyway.
			\item Use a competent specialist or decline.
			\item Ask management to perform it.
			\item Skip risky areas.
		\end{enumerate}
		\vspace{4pt}
		\textbf{Correct Answer: B}
	\end{frame}
	
	\begin{frame}[t]{Q47 -- Explanation}
		\small
		\textbf{Correct: B.} Standard 1210: auditors must possess or obtain the knowledge/skills needed, or use a competent specialist. \\[4pt]
		\textbf{A wrong:} Proceeding without competence violates the Standards. \\[2pt]
		\textbf{C wrong:} Management cannot replace internal audit. \\[2pt]
		\textbf{D wrong:} Skipping risky areas defeats the engagement's purpose.
	\end{frame}
	
	% ============ Q48 ============
	\begin{frame}[t]{Q48}
		\small
		\textbf{Question:} The internal audit plan should be based PRIMARILY on:
		\begin{enumerate}[label=\Alph*.,itemsep=1pt,topsep=2pt]
			\item Management requests.
			\item A risk assessment.
			\item Prior-year findings.
			\item Budget constraints.
		\end{enumerate}
		\vspace{4pt}
		\textbf{Correct Answer: B}
	\end{frame}
	
	\begin{frame}[t]{Q48 -- Explanation}
		\small
		\textbf{Correct: B.} Standard 2010: the plan must be based on a documented risk assessment, updated at least annually. \\[4pt]
		\textbf{A wrong:} Management input is considered but not primary. \\[2pt]
		\textbf{C wrong:} Prior findings inform, but risk drives. \\[2pt]
		\textbf{D wrong:} Budget is a constraint, not the basis.
	\end{frame}
	
	% ============ Q49 ============
	\begin{frame}[t]{Q49}
		\small
		\textbf{Question:} Which is the best description of "professional skepticism"?
		\begin{enumerate}[label=\Alph*.,itemsep=1pt,topsep=2pt]
			\item Assuming management is dishonest.
			\item A questioning mind and critical assessment of evidence.
			\item Ignoring evidence.
			\item Relying on prior-year conclusions.
		\end{enumerate}
		\vspace{4pt}
		\textbf{Correct Answer: B}
	\end{frame}
	
	\begin{frame}[t]{Q49 -- Explanation}
		\small
		\textbf{Correct: B.} Professional skepticism is a questioning mind, alertness to conditions suggesting misstatement, and critical evaluation of evidence. \\[4pt]
		\textbf{A wrong:} It is not cynicism or assuming dishonesty. \\[2pt]
		\textbf{C wrong:} Ignoring evidence is the opposite. \\[2pt]
		\textbf{D wrong:} Prior conclusions must be re-validated.
	\end{frame}
	
	% ============ Q50 ============
	\begin{frame}[t]{Q50}
		\small
		\textbf{Question:} Internal audit's role in ERM is to:
		\begin{enumerate}[label=\Alph*.,itemsep=1pt,topsep=2pt]
			\item Own the ERM process.
			\item Provide assurance and consulting on ERM.
			\item Set risk appetite.
			\item Approve risk responses.
		\end{enumerate}
		\vspace{4pt}
		\textbf{Correct Answer: B}
	\end{frame}
	
	\begin{frame}[t]{Q50 -- Explanation}
		\small
		\textbf{Correct: B.} Internal audit provides assurance and consulting on ERM but should not own it or set appetite -- that would impair objectivity. \\[4pt]
		\textbf{A wrong:} Management owns ERM. \\[2pt]
		\textbf{C wrong:} The board sets risk appetite. \\[2pt]
		\textbf{D wrong:} Management approves responses.
	\end{frame}
	
	% ============ Q51 ============
	\begin{frame}[t]{Q51}
		\small
		\textbf{Question:} Which is NOT a component of the IIA Code of Ethics?
		\begin{enumerate}[label=\Alph*.,itemsep=1pt,topsep=2pt]
			\item Integrity.
			\item Objectivity.
			\item Profitability.
			\item Confidentiality.
		\end{enumerate}
		\vspace{4pt}
		\textbf{Correct Answer: C}
	\end{frame}
	
	\begin{frame}[t]{Q51 -- Explanation}
		\small
		\textbf{Correct: C.} The Code's principles: Integrity, Objectivity, Confidentiality, Competency. Profitability is not one. \\[4pt]
		\textbf{A wrong:} Integrity is required. \\[2pt]
		\textbf{B wrong:} Objectivity is required. \\[2pt]
		\textbf{D wrong:} Confidentiality is required.
	\end{frame}
	
	% ============ Q52 ============
	\begin{frame}[t]{Q52}
		\small
		\textbf{Question:} Which is the most reliable evidence?
		\begin{enumerate}[label=\Alph*.,itemsep=1pt,topsep=2pt]
			\item Inquiry of management.
			\item Internally generated memos.
			\item External confirmations.
			\item Analytical review.
		\end{enumerate}
		\vspace{4pt}
		\textbf{Correct Answer: C}
	\end{frame}
	
	\begin{frame}[t]{Q52 -- Explanation}
		\small
		\textbf{Correct: C.} External confirmations come from independent third parties, making them highly reliable. \\[4pt]
		\textbf{A wrong:} Inquiry is weak. \\[2pt]
		\textbf{B wrong:} Internal memos are less reliable. \\[2pt]
		\textbf{D wrong:} Analytical review is indirect.
	\end{frame}
	
	% ============ Q53 ============
	\begin{frame}[t]{Q53}
		\small
		\textbf{Question:} Which is a limitation of internal control?
		\begin{enumerate}[label=\Alph*.,itemsep=1pt,topsep=2pt]
			\item Cost-benefit considerations.
			\item Segregation of duties.
			\item Physical safeguards.
			\item Authorization procedures.
		\end{enumerate}
		\vspace{4pt}
		\textbf{Correct Answer: A}
	\end{frame}
	
	\begin{frame}[t]{Q53 -- Explanation}
		\small
		\textbf{Correct: A.} Controls are limited by cost-benefit, human error, collusion, and management override. \\[4pt]
		\textbf{B wrong:} Segregation of duties is a control, not a limitation. \\[2pt]
		\textbf{C wrong:} Safeguards are controls. \\[2pt]
		\textbf{D wrong:} Authorization is a control.
	\end{frame}
	
	% ============ Q54 ============
	\begin{frame}[t]{Q54}
		\small
		\textbf{Question:} Auditors should document:
		\begin{enumerate}[label=\Alph*.,itemsep=1pt,topsep=2pt]
			\item Only significant findings.
			\item Only final conclusions.
			\item The work performed and conclusions reached.
			\item Only what management requests.
		\end{enumerate}
		\vspace{4pt}
		\textbf{Correct Answer: C}
	\end{frame}
	
	\begin{frame}[t]{Q54 -- Explanation}
		\small
		\textbf{Correct: C.} Standard 2330: workpapers must document the work performed, evidence obtained, and conclusions reached. \\[4pt]
		\textbf{A wrong:} All relevant work is documented. \\[2pt]
		\textbf{B wrong:} Conclusions alone lack support. \\[2pt]
		\textbf{D wrong:} Documentation is not management's choice.
	\end{frame}
	
	% ============ Q55 ============
	\begin{frame}[t]{Q55}
		\small
		\textbf{Question:} The best way to test IT general controls is to:
		\begin{enumerate}[label=\Alph*.,itemsep=1pt,topsep=2pt]
			\item Rely on management's assertion.
			\item Test access, change management, and operations controls.
			\item Review the IT budget.
			\item Interview the CIO only.
		\end{enumerate}
		\vspace{4pt}
		\textbf{Correct Answer: B}
	\end{frame}
	
	\begin{frame}[t]{Q55 -- Explanation}
		\small
		\textbf{Correct: B.} ITGCs cover access to programs/data, program changes, and computer operations. Testing these gives assurance over the IT environment. \\[4pt]
		\textbf{A wrong:} Assertions alone are insufficient. \\[2pt]
		\textbf{C wrong:} Budget review is not a control test. \\[2pt]
		\textbf{D wrong:} Interviews alone are weak evidence.
	\end{frame}
	
	% ============ Q56 ============
	\begin{frame}[t]{Q56}
		\small
		\textbf{Question:} An auditor disagrees with management's response to a finding. The auditor should:
		\begin{enumerate}[label=\Alph*.,itemsep=1pt,topsep=2pt]
			\item Accept management's response.
			\item Document the disagreement and escalate if needed.
			\item Ignore the finding.
			\item Resign.
		\end{enumerate}
		\vspace{4pt}
		\textbf{Correct Answer: B}
	\end{frame}
	
	\begin{frame}[t]{Q56 -- Explanation}
		\small
		\textbf{Correct: B.} The auditor documents management's response and their own disagreement, escalating to the board if the issue is significant. \\[4pt]
		\textbf{A wrong:} Accepting without challenge undermines audit. \\[2pt]
		\textbf{C wrong:} Ignoring findings violates the Standards. \\[2pt]
		\textbf{D wrong:} Resignation is a last resort.
	\end{frame}
	
	% ============ Q57 ============
	\begin{frame}[t]{Q57}
		\small
		\textbf{Question:} The scope of an engagement should be:
		\begin{enumerate}[label=\Alph*.,itemsep=1pt,topsep=2pt]
			\item Determined by management.
			\item Sufficient to achieve engagement objectives.
			\item Limited to financial areas.
			\item Limited to one year.
		\end{enumerate}
		\vspace{4pt}
		\textbf{Correct Answer: B}
	\end{frame}
	
	\begin{frame}[t]{Q57 -- Explanation}
		\small
		\textbf{Correct: B.} Scope must be sufficient to meet the engagement's objectives, as determined during planning. \\[4pt]
		\textbf{A wrong:} Management may request, but scope is internal audit's professional judgment. \\[2pt]
		\textbf{C wrong:} Scope can cover operations, compliance, IT, etc. \\[2pt]
		\textbf{D wrong:} No fixed time limit.
	\end{frame}
	
	% ============ Q58 ============
	\begin{frame}[t]{Q58}
		\small
		\textbf{Question:} Which is the best evidence of the existence of accounts receivable?
		\begin{enumerate}[label=\Alph*.,itemsep=1pt,topsep=2pt]
			\item Inquiry of the credit manager.
			\item External confirmation from customers.
			\item The aged trial balance.
			\item The sales budget.
		\end{enumerate}
		\vspace{4pt}
		\textbf{Correct Answer: B}
	\end{frame}
	
	\begin{frame}[t]{Q58 -- Explanation}
		\small
		\textbf{Correct: B.} External confirmations from customers directly verify existence and are highly reliable. \\[4pt]
		\textbf{A wrong:} Inquiry is weak. \\[2pt]
		\textbf{C wrong:} The trial balance is internally generated. \\[2pt]
		\textbf{D wrong:} Budgets are forward-looking, not evidence of existence.
	\end{frame}
	
	% ============ Q59 ============
	\begin{frame}[t]{Q59}
		\small
		\textbf{Question:} The CAE should ensure internal auditors:
		\begin{enumerate}[label=\Alph*.,itemsep=1pt,topsep=2pt]
			\item Never make mistakes.
			\item Have the knowledge, skills, and competencies needed.
			\item Only hold accounting degrees.
			\item Work alone.
		\end{enumerate}
		\vspace{4pt}
		\textbf{Correct Answer: B}
	\end{frame}
	
	\begin{frame}[t]{Q59 -- Explanation}
		\small
		\textbf{Correct: B.} Standard 1210: auditors must possess the knowledge, skills, and other competencies needed to perform their responsibilities. \\[4pt]
		\textbf{A wrong:} No one is error-free. \\[2pt]
		\textbf{C wrong:} Diverse skills are needed, not only accounting. \\[2pt]
		\textbf{D wrong:} Teamwork and supervision are expected.
	\end{frame}
	
	% ============ Q60 ============
	\begin{frame}[t]{Q60}
		\small
		\textbf{Question:} Which best describes control self-assessment (CSA)?
		\begin{enumerate}[label=\Alph*.,itemsep=1pt,topsep=2pt]
			\item An external audit technique.
			\item A workshop-based approach where management assesses controls.
			\item A type of fraud.
			\item A regulatory filing.
		\end{enumerate}
		\vspace{4pt}
		\textbf{Correct Answer: B}
	\end{frame}
	
	\begin{frame}[t]{Q60 -- Explanation}
		\small
		\textbf{Correct: B.} CSA engages management/process owners in assessing their own controls, often in facilitated workshops. Internal audit may facilitate. \\[4pt]
		\textbf{A wrong:} CSA is internal. \\[2pt]
		\textbf{C wrong:} It is not fraud. \\[2pt]
		\textbf{D wrong:} It is not a filing.
	\end{frame}
	
	% ============ Q61 ============
	\begin{frame}[t]{Q61}
		\small
		\textbf{Question:} Which is the best way to obtain evidence about the design of a control?
		\begin{enumerate}[label=\Alph*.,itemsep=1pt,topsep=2pt]
			\item Reperformance.
			\item Inquiry and inspection of documents.
			\item Observation only.
			\item Analytical review.
		\end{enumerate}
		\vspace{4pt}
		\textbf{Correct Answer: B}
	\end{frame}
	
	\begin{frame}[t]{Q61 -- Explanation}
		\small
		\textbf{Correct: B.} Design testing uses inquiry and inspection of policies/procedures to see if the control is properly designed. \\[4pt]
		\textbf{A wrong:} Reperformance tests operating effectiveness. \\[2pt]
		\textbf{C wrong:} Observation alone is insufficient. \\[2pt]
		\textbf{D wrong:} Analytical review does not test design.
	\end{frame}
	
	% ============ Q62 ============
	\begin{frame}[t]{Q62}
		\small
		\textbf{Question:} The CAE reports administratively to the CEO. This means the CEO:
		\begin{enumerate}[label=\Alph*.,itemsep=1pt,topsep=2pt]
			\item Approves the audit charter.
			\item Handles HR, budget, and day-to-day administration.
			\item Sets the audit plan.
			\item Issues audit opinions.
		\end{enumerate}
		\vspace{4pt}
		\textbf{Correct Answer: B}
	\end{frame}
	
	\begin{frame}[t]{Q62 -- Explanation}
		\small
		\textbf{Correct: B.} Administrative reporting covers HR, budgets, and internal communication, while functional reporting (to the board) covers charter, plan, and reporting. \\[4pt]
		\textbf{A wrong:} The board approves the charter. \\[2pt]
		\textbf{C wrong:} The board approves the risk-based plan. \\[2pt]
		\textbf{D wrong:} Internal audit doesn't issue formal opinions.
	\end{frame}
	
	% ============ Q63 ============
	\begin{frame}[t]{Q63}
		\small
		\textbf{Question:} Which is an indicator of potential fraud?
		\begin{enumerate}[label=\Alph*.,itemsep=1pt,topsep=2pt]
			\item Strong segregation of duties.
			\item Lifestyle changes inconsistent with salary.
			\item Regular reconciliations.
			\item Timely reporting.
		\end{enumerate}
		\vspace{4pt}
		\textbf{Correct Answer: B}
	\end{frame}
	
	\begin{frame}[t]{Q63 -- Explanation}
		\small
		\textbf{Correct: B.} Unexplained lifestyle changes are a classic red flag for fraud. \\[4pt]
		\textbf{A wrong:} Segregation of duties is a control. \\[2pt]
		\textbf{C wrong:} Reconciliations are controls. \\[2pt]
		\textbf{D wrong:} Timely reporting is a positive sign.
	\end{frame}
	
	% ============ Q64 ============
	\begin{frame}[t]{Q64}
		\small
		\textbf{Question:} Which is the best way to maintain independence when consulting?
		\begin{enumerate}[label=\Alph*.,itemsep=1pt,topsep=2pt]
			\item Avoid documenting the engagement.
			\item Ensure the auditor does not assume management responsibilities.
			\item Have the auditor implement the recommendations.
			\item Skip the engagement letter.
		\end{enumerate}
		\vspace{4pt}
		\textbf{Correct Answer: B}
	\end{frame}
	
	\begin{frame}[t]{Q64 -- Explanation}
		\small
		\textbf{Correct: B.} Internal auditors may consult but must not take on management responsibilities, which would impair objectivity. \\[4pt]
		\textbf{A wrong:} All engagements must be documented. \\[2pt]
		\textbf{C wrong:} Implementing decisions is a management function. \\[2pt]
		\textbf{D wrong:} Engagement letters are good practice.
	\end{frame}
	
	% ============ Q65 ============
	\begin{frame}[t]{Q65}
		\small
		\textbf{Question:} The primary reason for a risk-based audit plan is to:
		\begin{enumerate}[label=\Alph*.,itemsep=1pt,topsep=2pt]
			\item Reduce cost.
			\item Focus on areas of greatest risk.
			\item Speed up fieldwork.
			\item Reduce findings.
		\end{enumerate}
		\vspace{4pt}
		\textbf{Correct Answer: B}
	\end{frame}
	
	\begin{frame}[t]{Q65 -- Explanation}
		\small
		\textbf{Correct: B.} Risk-based planning directs limited audit resources to the highest-risk areas, maximizing value. \\[4pt]
		\textbf{A wrong:} Cost is not the primary driver. \\[2pt]
		\textbf{C wrong:} Speed is not the objective. \\[2pt]
		\textbf{D wrong:} Fewer findings is not the goal.
	\end{frame}
	
	% ============ Q66 ============
	\begin{frame}[t]{Q66}
		\small
		\textbf{Question:} Which is the best description of "reasonable assurance"?
		\begin{enumerate}[label=\Alph*.,itemsep=1pt,topsep=2pt]
			\item A guarantee of no errors.
			\item A high but not absolute level of assurance.
			\item No assurance at all.
			\item 100\% certainty.
		\end{enumerate}
		\vspace{4pt}
		\textbf{Correct Answer: B}
	\end{frame}
	
	\begin{frame}[t]{Q66 -- Explanation}
		\small
		\textbf{Correct: B.} Reasonable assurance is a high level of assurance, but not absolute, due to inherent limitations of audit. \\[4pt]
		\textbf{A wrong:} No guarantee is possible. \\[2pt]
		\textbf{C wrong:} It is not zero assurance. \\[2pt]
		\textbf{D wrong:} Absolute certainty is unattainable.
	\end{frame}
	
	% ============ Q67 ============
	\begin{frame}[t]{Q67}
		\small
		\textbf{Question:} Which document authorizes internal audit's access to records and personnel?
		\begin{enumerate}[label=\Alph*.,itemsep=1pt,topsep=2pt]
			\item The engagement letter.
			\item The internal audit charter.
			\item The audit plan.
			\item The Code of Ethics.
		\end{enumerate}
		\vspace{4pt}
		\textbf{Correct Answer: B}
	\end{frame}
	
	\begin{frame}[t]{Q67 -- Explanation}
		\small
		\textbf{Correct: B.} The charter grants internal audit authority, including access to records, personnel, and physical properties. \\[4pt]
		\textbf{A wrong:} Engagement letters are for specific engagements. \\[2pt]
		\textbf{C wrong:} The plan lists engagements, not authority. \\[2pt]
		\textbf{D wrong:} The Code guides conduct, not authority.
	\end{frame}
	
	% ============ Q68 ============
	\begin{frame}[t]{Q68}
		\small
		\textbf{Question:} The CAE should report to the board if:
		\begin{enumerate}[label=\Alph*.,itemsep=1pt,topsep=2pt]
			\item Management disagrees with a minor observation.
			\item There is a significant risk or control issue management will not address.
			\item The audit budget is exceeded.
			\item An auditor is late.
		\end{enumerate}
		\vspace{4pt}
		\textbf{Correct Answer: B}
	\end{frame}
	
	\begin{frame}[t]{Q68 -- Explanation}
		\small
		\textbf{Correct: B.} Significant unresolved risk/control issues must be escalated to the board. \\[4pt]
		\textbf{A wrong:} Minor disagreements are handled at the working level. \\[2pt]
		\textbf{C wrong:} Budget variances are administrative. \\[2pt]
		\textbf{D wrong:} Tardiness is an HR matter.
	\end{frame}
	
	% ============ Q69 ============
	\begin{frame}[t]{Q69}
		\small
		\textbf{Question:} Which is the best way to test the completeness of revenue?
		\begin{enumerate}[label=\Alph*.,itemsep=1pt,topsep=2pt]
			\item Confirm accounts receivable.
			\item Trace from shipping documents to sales invoices.
			\item Review the sales budget.
			\item Interview the sales manager.
		\end{enumerate}
		\vspace{4pt}
		\textbf{Correct Answer: B}
	\end{frame}
	
	\begin{frame}[t]{Q69 -- Explanation}
		\small
		\textbf{Correct: B.} Tracing from shipping documents (source) to invoices tests completeness -- did all shipments get billed? \\[4pt]
		\textbf{A wrong:} Confirmations test existence, not completeness. \\[2pt]
		\textbf{C wrong:} Budgets are not evidence. \\[2pt]
		\textbf{D wrong:} Inquiry alone is weak.
	\end{frame}
	
	% ============ Q70 ============
	\begin{frame}[t]{Q70}
		\small
		\textbf{Question:} Which is the best way to test the existence of revenue?
		\begin{enumerate}[label=\Alph*.,itemsep=1pt,topsep=2pt]
			\item Trace from sales invoices to shipping documents.
			\item Trace from shipping documents to sales invoices.
			\item Review the sales budget.
			\item Interview the sales manager.
		\end{enumerate}
		\vspace{4pt}
		\textbf{Correct Answer: A}
	\end{frame}
	
	\begin{frame}[t]{Q70 -- Explanation}
		\small
		\textbf{Correct: A.} Tracing from recorded invoices back to shipping documents tests existence -- did the sale really occur? \\[4pt]
		\textbf{B wrong:} That direction tests completeness. \\[2pt]
		\textbf{C wrong:} Budgets are not evidence. \\[2pt]
		\textbf{D wrong:} Inquiry alone is insufficient.
	\end{frame}
	
	% ============ Q71 ============
	\begin{frame}[t]{Q71}
		\small
		\textbf{Question:} Which is NOT a type of assurance engagement?
		\begin{enumerate}[label=\Alph*.,itemsep=1pt,topsep=2pt]
			\item Financial audit.
			\item Compliance audit.
			\item Operational audit.
			\item Advisory on system design.
		\end{enumerate}
		\vspace{4pt}
		\textbf{Correct Answer: D}
	\end{frame}
	
	\begin{frame}[t]{Q71 -- Explanation}
		\small
		\textbf{Correct: D.} Advisory on system design is consulting, not assurance. \\[4pt]
		\textbf{A wrong:} Financial audit is assurance. \\[2pt]
		\textbf{B wrong:} Compliance audit is assurance. \\[2pt]
		\textbf{C wrong:} Operational audit is assurance.
	\end{frame}
	
	% ============ Q72 ============
	\begin{frame}[t]{Q72}
		\small
		\textbf{Question:} An internal auditor should be most concerned with which of the following?
		\begin{enumerate}[label=\Alph*.,itemsep=1pt,topsep=2pt]
			\item A control that is documented but not tested.
			\item A control that is designed but not operating.
			\item A control that is operating but not documented.
			\item A control that is both designed and operating.
		\end{enumerate}
		\vspace{4pt}
		\textbf{Correct Answer: B}
	\end{frame}
	
	\begin{frame}[t]{Q72 -- Explanation}
		\small
		\textbf{Correct: B.} A control that exists on paper but does not operate provides no protection -- a serious deficiency. \\[4pt]
		\textbf{A wrong:} Documentation without testing is a scope issue, less severe. \\[2pt]
		\textbf{C wrong:} Operating but undocumented is a documentation gap. \\[2pt]
		\textbf{D wrong:} Designed and operating is ideal.
	\end{frame}
	
	% ============ Q73 ============
	\begin{frame}[t]{Q73}
		\small
		\textbf{Question:} Which is the best description of "governance"?
		\begin{enumerate}[label=\Alph*.,itemsep=1pt,topsep=2pt]
			\item The process of preparing budgets.
			\item The processes and structures to direct and control the organization.
			\item The audit committee's meeting minutes.
			\item The external auditor's report.
		\end{enumerate}
		\vspace{4pt}
		\textbf{Correct Answer: B}
	\end{frame}
	
	\begin{frame}[t]{Q73 -- Explanation}
		\small
		\textbf{Correct: B.} Governance encompasses the processes and structures used to direct, control, and hold the organization accountable. \\[4pt]
		\textbf{A wrong:} Budgeting is a management process. \\[2pt]
		\textbf{C wrong:} Minutes are evidence, not governance. \\[2pt]
		\textbf{D wrong:} The audit report is an output.
	\end{frame}
	
	% ============ Q74 ============
	\begin{frame}[t]{Q74}
		\small
		\textbf{Question:} The board's role in governance includes:
		\begin{enumerate}[label=\Alph*.,itemsep=1pt,topsep=2pt]
			\item Daily operational decisions.
			\item Oversight of risk, control, and accountability to stakeholders.
			\item Preparing workpapers.
			\item Hiring all staff.
		\end{enumerate}
		\vspace{4pt}
		\textbf{Correct Answer: B}
	\end{frame}
	
	\begin{frame}[t]{Q74 -- Explanation}
		\small
		\textbf{Correct: B.} The board oversees management, risk, and control, and is accountable to stakeholders. \\[4pt]
		\textbf{A wrong:} Daily operations are management's job. \\[2pt]
		\textbf{C wrong:} Workpapers are internal audit's. \\[2pt]
		\textbf{D wrong:} Hiring is management/HR.
	\end{frame}
	
	% ============ Q75 ============
	\begin{frame}[t]{Q75}
		\small
		\textbf{Question:} Which is the best description of "risk response"?
		\begin{enumerate}[label=\Alph*.,itemsep=1pt,topsep=2pt]
			\item Ignoring risk.
			\item Avoid, reduce, share, or accept risk.
			\item Transferring all risk.
			\item Eliminating all risk.
		\end{enumerate}
		\vspace{4pt}
		\textbf{Correct Answer: B}
	\end{frame}
	
	\begin{frame}[t]{Q75 -- Explanation}
		\small
		\textbf{Correct: B.} The four responses are avoid, reduce (mitigate), share (transfer), and accept. \\[4pt]
		\textbf{A wrong:} Ignoring is not a valid response. \\[2pt]
		\textbf{C wrong:} Transfer is only one option. \\[2pt]
		\textbf{D wrong:} Elimination is not always possible.
	\end{frame}
	
	% ============ Q76 ============
	\begin{frame}[t]{Q76}
		\small
		\textbf{Question:} Which is the best description of "control activities"?
		\begin{enumerate}[label=\Alph*.,itemsep=1pt,topsep=2pt]
			\item Policies and procedures that help ensure management directives are carried out.
			\item The board's meeting schedule.
			\item The audit plan.
			\item The marketing plan.
		\end{enumerate}
		\vspace{4pt}
		\textbf{Correct Answer: A}
	\end{frame}
	
	\begin{frame}[t]{Q76 -- Explanation}
		\small
		\textbf{Correct: A.} Control activities are the policies and procedures (approvals, reconciliations, segregation) that help ensure directives are executed. \\[4pt]
		\textbf{B wrong:} Meeting schedules are not control activities. \\[2pt]
		\textbf{C wrong:} The audit plan is a planning document. \\[2pt]
		\textbf{D wrong:} Marketing plans are unrelated.
	\end{frame}
	
	% ============ Q77 ============
	\begin{frame}[t]{Q77}
		\small
		\textbf{Question:} Which is the best description of "monitoring" in COSO?
		\begin{enumerate}[label=\Alph*.,itemsep=1pt,topsep=2pt]
			\item Preparing the budget.
			\item Assessing the quality of internal control over time.
			\item Hiring staff.
			\item Setting the risk appetite.
		\end{enumerate}
		\vspace{4pt}
		\textbf{Correct Answer: B}
	\end{frame}
	
	\begin{frame}[t]{Q77 -- Explanation}
		\small
		\textbf{Correct: B.} Monitoring assesses the quality of internal control performance over time, through ongoing and separate evaluations. \\[4pt]
		\textbf{A wrong:} Budgeting is a control activity. \\[2pt]
		\textbf{C wrong:} Hiring is a control activity. \\[2pt]
		\textbf{D wrong:} Risk appetite is set by the board.
	\end{frame}
	
	% ============ Q78 ============
	\begin{frame}[t]{Q78}
		\small
		\textbf{Question:} Which is the best description of "information and communication" in COSO?
		\begin{enumerate}[label=\Alph*.,itemsep=1pt,topsep=2pt]
			\item The marketing campaign.
			\item Identifying, capturing, and communicating relevant information.
			\item The IT budget.
			\item The audit report.
		\end{enumerate}
		\vspace{4pt}
		\textbf{Correct Answer: B}
	\end{frame}
	
	\begin{frame}[t]{Q78 -- Explanation}
		\small
		\textbf{Correct: B.} Information and communication ensures relevant information is identified, captured, and communicated in a timely manner. \\[4pt]
		\textbf{A wrong:} Marketing is unrelated. \\[2pt]
		\textbf{C wrong:} IT budget is a resource, not a COSO component. \\[2pt]
		\textbf{D wrong:} The audit report is an output.
	\end{frame}
	
	% ============ Q79 ============
	\begin{frame}[t]{Q79}
		\small
		\textbf{Question:} Which is the best description of "risk assessment" in COSO?
		\begin{enumerate}[label=\Alph*.,itemsep=1pt,topsep=2pt]
			\item Identifying and analyzing risks to objectives.
			\item Preparing financial statements.
			\item Hiring staff.
			\item Auditing the financial statements.
		\end{enumerate}
		\vspace{4pt}
		\textbf{Correct Answer: A}
	\end{frame}
	
	\begin{frame}[t]{Q79 -- Explanation}
		\small
		\textbf{Correct: A.} Risk assessment identifies and analyzes risks to achieving objectives, forming a basis for how risks should be managed. \\[4pt]
		\textbf{B wrong:} Financial statement prep is a management function. \\[2pt]
		\textbf{C wrong:} Hiring is a control activity. \\[2pt]
		\textbf{D wrong:} Auditing is assurance, not risk assessment.
	\end{frame}
	
	% ============ Q80 ============
	\begin{frame}[t]{Q80}
		\small
		\textbf{Question:} The control environment is best described as:
		\begin{enumerate}[label=\Alph*.,itemsep=1pt,topsep=2pt]
			\item The tone of the organization, influencing control consciousness.
			\item The IT infrastructure.
			\item The audit committee.
			\item The external auditor's opinion.
		\end{enumerate}
		\vspace{4pt}
		\textbf{Correct Answer: A}
	\end{frame}
	
	\begin{frame}[t]{Q80 -- Explanation}
		\small
		\textbf{Correct: A.} The control environment is the tone at the top -- integrity, ethical values, and commitment to competence -- that sets the foundation for all other controls. \\[4pt]
		\textbf{B wrong:} IT infrastructure is part of control activities. \\[2pt]
		\textbf{C wrong:} The audit committee is part of governance, not the whole environment. \\[2pt]
		\textbf{D wrong:} The external opinion is not the environment.
	\end{frame}
	
	% ============ Q81 ============
	\begin{frame}[t]{Q81}
		\small
		\textbf{Question:} Which is the best description of "internal control"?
		\begin{enumerate}[label=\Alph*.,itemsep=1pt,topsep=2pt]
			\item A process designed to provide reasonable assurance regarding objectives.
			\item A guarantee of no fraud.
			\item The external audit opinion.
			\item The marketing plan.
		\end{enumerate}
		\vspace{4pt}
		\textbf{Correct Answer: A}
	\end{frame}
	
	\begin{frame}[t]{Q81 -- Explanation}
		\small
		\textbf{Correct: A.} Internal control is a process, effected by people, designed to provide reasonable assurance regarding operations, reporting, and compliance objectives. \\[4pt]
		\textbf{B wrong:} No control guarantees no fraud. \\[2pt]
		\textbf{C wrong:} The external opinion is separate. \\[2pt]
		\textbf{D wrong:} Marketing is unrelated.
	\end{frame}
	
	% ============ Q82 ============
	\begin{frame}[t]{Q82}
		\small
		\textbf{Question:} Which is the best description of "compliance risk"?
		\begin{enumerate}[label=\Alph*.,itemsep=1pt,topsep=2pt]
			\item Risk of violating laws, regulations, or policies.
			\item Risk of losing market share.
			\item Risk of IT failure.
			\item Risk of fraud.
		\end{enumerate}
		\vspace{4pt}
		\textbf{Correct Answer: A}
	\end{frame}
	
	\begin{frame}[t]{Q82 -- Explanation}
		\small
		\textbf{Correct: A.} Compliance risk is the risk of violating laws, regulations, rules, or policies, leading to penalties or reputational damage. \\[4pt]
		\textbf{B wrong:} Market share is strategic risk. \\[2pt]
		\textbf{C wrong:} IT failure is operational/technology risk. \\[2pt]
		\textbf{D wrong:} Fraud is its own category.
	\end{frame}
	
	% ============ Q83 ============
	\begin{frame}[t]{Q83}
		\small
		\textbf{Question:} Which is the best description of "operational risk"?
		\begin{enumerate}[label=\Alph*.,itemsep=1pt,topsep=2pt]
			\item Risk of loss from inadequate or failed processes, people, or systems.
			\item Risk of losing market share.
			\item Risk of violating laws.
			\item Risk of fraud.
		\end{enumerate}
		\vspace{4pt}
		\textbf{Correct Answer: A}
	\end{frame}
	
	\begin{frame}[t]{Q83 -- Explanation}
		\small
		\textbf{Correct: A.} Operational risk arises from inadequate or failed internal processes, people, systems, or external events. \\[4pt]
		\textbf{B wrong:} Market share is strategic. \\[2pt]
		\textbf{C wrong:} Legal violations are compliance risk. \\[2pt]
		\textbf{D wrong:} Fraud is a subset, not the definition.
	\end{frame}
	
	% ============ Q84 ============
	\begin{frame}[t]{Q84}
		\small
		\textbf{Question:} Which is the best description of "strategic risk"?
		\begin{enumerate}[label=\Alph*.,itemsep=1pt,topsep=2pt]
			\item Risk of failing to achieve strategic objectives.
			\item Risk of IT failure.
			\item Risk of fraud.
			\item Risk of violating laws.
		\end{enumerate}
		\vspace{4pt}
		\textbf{Correct Answer: A}
	\end{frame}
	
	\begin{frame}[t]{Q84 -- Explanation}
		\small
		\textbf{Correct: A.} Strategic risk is the risk of failing to achieve the organization's strategic objectives due to poor decisions or external changes. \\[4pt]
		\textbf{B wrong:} IT failure is operational. \\[2pt]
		\textbf{C wrong:} Fraud is its own category. \\[2pt]
		\textbf{D wrong:} Legal violations are compliance.
	\end{frame}
	
	% ============ Q85 ============
	\begin{frame}[t]{Q85}
		\small
		\textbf{Question:} Which is the best description of "reputational risk"?
		\begin{enumerate}[label=\Alph*.,itemsep=1pt,topsep=2pt]
			\item Risk of damage to the organization's reputation.
			\item Risk of IT failure.
			\item Risk of fraud.
			\item Risk of violating laws.
		\end{enumerate}
		\vspace{4pt}
		\textbf{Correct Answer: A}
	\end{frame}
	
	\begin{frame}[t]{Q85 -- Explanation}
		\small
		\textbf{Correct: A.} Reputational risk is the risk of damage to the organization's image, brand, or public trust. \\[4pt]
		\textbf{B wrong:} IT failure is operational. \\[2pt]
		\textbf{C wrong:} Fraud is a separate category. \\[2pt]
		\textbf{D wrong:} Legal violations are compliance risk.
	\end{frame}
	
	% ============ Q86 ============
	\begin{frame}[t]{Q86}
		\small
		\textbf{Question:} Which is the best description of "risk register"?
		\begin{enumerate}[label=\Alph*.,itemsep=1pt,topsep=2pt]
			\item A list of identified risks and their responses.
			\item The audit plan.
			\item The marketing plan.
			\item The financial statements.
		\end{enumerate}
		\vspace{4pt}
		\textbf{Correct Answer: A}
	\end{frame}
	
	\begin{frame}[t]{Q86 -- Explanation}
		\small
		\textbf{Correct: A.} A risk register documents identified risks, their assessment, owners, and responses. \\[4pt]
		\textbf{B wrong:} The audit plan is separate. \\[2pt]
		\textbf{C wrong:} The marketing plan is unrelated. \\[2pt]
		\textbf{D wrong:} Financial statements are outputs.
	\end{frame}
	
	% ============ Q87 ============
	\begin{frame}[t]{Q87}
		\small
		\textbf{Question:} Which is the best description of "control risk"?
		\begin{enumerate}[label=\Alph*.,itemsep=1pt,topsep=2pt]
			\item The risk that a control fails to prevent or detect a misstatement.
			\item The risk of fraud.
			\item The risk of IT failure.
			\item The risk of market loss.
		\end{enumerate}
		\vspace{4pt}
		\textbf{Correct Answer: A}
	\end{frame}
	
	\begin{frame}[t]{Q87 -- Explanation}
		\small
		\textbf{Correct: A.} Control risk is the risk that a control will not prevent or detect a material misstatement on a timely basis. \\[4pt]
		\textbf{B wrong:} Fraud is a source, not control risk. \\[2pt]
		\textbf{C wrong:} IT failure is operational. \\[2pt]
		\textbf{D wrong:} Market loss is strategic/market risk.
	\end{frame}
	
	% ============ Q88 ============
	\begin{frame}[t]{Q88}
		\small
		\textbf{Question:} Which is the best description of "detection risk"?
		\begin{enumerate}[label=\Alph*.,itemsep=1pt,topsep=2pt]
			\item The risk that audit procedures fail to detect a misstatement.
			\item The risk of fraud.
			\item The risk of IT failure.
			\item The risk of market loss.
		\end{enumerate}
		\vspace{4pt}
		\textbf{Correct Answer: A}
	\end{frame}
	
	\begin{frame}[t]{Q88 -- Explanation}
		\small
		\textbf{Correct: A.} Detection risk is the risk that the auditor's procedures will not detect a material misstatement. It is controlled by the auditor. \\[4pt]
		\textbf{B wrong:} Fraud is separate. \\[2pt]
		\textbf{C wrong:} IT failure is operational. \\[2pt]
		\textbf{D wrong:} Market loss is market risk.
	\end{frame}
	
	% ============ Q89 ============
	\begin{frame}[t]{Q89}
		\small
		\textbf{Question:} Which is the best description of "audit evidence"?
		\begin{enumerate}[label=\Alph*.,itemsep=1pt,topsep=2pt]
			\item Information used by the auditor to reach conclusions.
			\item The audit report.
			\item The audit plan.
			\item The engagement letter.
		\end{enumerate}
		\vspace{4pt}
		\textbf{Correct Answer: A}
	\end{frame}
	
	\begin{frame}[t]{Q89 -- Explanation}
		\small
		\textbf{Correct: A.} Audit evidence is all information used by the auditor to determine whether the engagement objectives are met and to reach conclusions. \\[4pt]
		\textbf{B wrong:} The report is the output. \\[2pt]
		\textbf{C wrong:} The plan is a planning document. \\[2pt]
		\textbf{D wrong:} The engagement letter defines the agreement.
	\end{frame}
	
	% ============ Q90 ============
	\begin{frame}[t]{Q90}
		\small
		\textbf{Question:} Which is the best description of "workpapers"?
		\begin{enumerate}[label=\Alph*.,itemsep=1pt,topsep=2pt]
			\item Records documenting the audit work performed and conclusions reached.
			\item The audit report.
			\item The audit plan.
			\item The engagement letter.
		\end{enumerate}
		\vspace{4pt}
		\textbf{Correct Answer: A}
	\end{frame}
	
	\begin{frame}[t]{Q90 -- Explanation}
		\small
		\textbf{Correct: A.} Workpapers document the work performed, evidence obtained, and conclusions reached, providing support for the report. \\[4pt]
		\textbf{B wrong:} The report is a separate deliverable. \\[2pt]
		\textbf{C wrong:} The plan is a planning document. \\[2pt]
		\textbf{D wrong:} The engagement letter defines the agreement.
	\end{frame}
	
	% ============ Q91 ============
	\begin{frame}[t]{Q91}
		\small
		\textbf{Question:} Which is the best description of "audit program"?
		\begin{enumerate}[label=\Alph*.,itemsep=1pt,topsep=2pt]
			\item A document listing the procedures to be performed.
			\item The audit report.
			\item The engagement letter.
			\item The audit committee minutes.
		\end{enumerate}
		\vspace{4pt}
		\textbf{Correct Answer: A}
	\end{frame}
	
	\begin{frame}[t]{Q91 -- Explanation}
		\small
		\textbf{Correct: A.} The audit program lists the procedures to be performed to achieve the engagement objectives, and is tailored to the engagement. \\[4pt]
		\textbf{B wrong:} The report is the output. \\[2pt]
		\textbf{C wrong:} The engagement letter defines the agreement. \\[2pt]
		\textbf{D wrong:} Minutes are governance records.
	\end{frame}
	
	% ============ Q92 ============
	\begin{frame}[t]{Q92}
		\small
		\textbf{Question:} Which is the best description of "engagement objectives"?
		\begin{enumerate}[label=\Alph*.,itemsep=1pt,topsep=2pt]
			\item What the engagement is intended to achieve.
			\item The audit fee.
			\item The audit committee's schedule.
			\item The marketing plan.
		\end{enumerate}
		\vspace{4pt}
		\textbf{Correct Answer: A}
	\end{frame}
	
	\begin{frame}[t]{Q92 -- Explanation}
		\small
		\textbf{Correct: A.} Engagement objectives state what the engagement is intended to accomplish, guiding scope and procedures. \\[4pt]
		\textbf{B wrong:} Fees are administrative. \\[2pt]
		\textbf{C wrong:} The committee's schedule is unrelated. \\[2pt]
		\textbf{D wrong:} Marketing is unrelated.
	\end{frame}
	
	% ============ Q93 ============
	\begin{frame}[t]{Q93}
		\small
		\textbf{Question:} Which is the best description of "engagement scope"?
		\begin{enumerate}[label=\Alph*.,itemsep=1pt,topsep=2pt]
			\item The boundaries of the engagement -- what is included and excluded.
			\item The audit fee.
			\item The audit committee's schedule.
			\item The marketing plan.
		\end{enumerate}
		\vspace{4pt}
		\textbf{Correct Answer: A}
	\end{frame}
	
	\begin{frame}[t]{Q93 -- Explanation}
		\small
		\textbf{Correct: A.} Scope defines the boundaries of the engagement -- the activities, periods, and locations covered. \\[4pt]
		\textbf{B wrong:} Fees are administrative. \\[2pt]
		\textbf{C wrong:} The committee's schedule is unrelated. \\[2pt]
		\textbf{D wrong:} Marketing is unrelated.
	\end{frame}
	
	% ============ Q94 ============
	\begin{frame}[t]{Q94}
		\small
		\textbf{Question:} Which is the best description of "engagement risk"?
		\begin{enumerate}[label=\Alph*.,itemsep=1pt,topsep=2pt]
			\item The risk that the engagement will fail to achieve its objectives or cause harm.
			\item The risk of fraud.
			\item The risk of IT failure.
			\item The risk of market loss.
		\end{enumerate}
		\vspace{4pt}
		\textbf{Correct Answer: A}
	\end{frame}
	
	\begin{frame}[t]{Q94 -- Explanation}
		\small
		\textbf{Correct: A.} Engagement risk is the risk that the engagement will not achieve its objectives or will cause harm to the organization or the auditor. \\[4pt]
		\textbf{B wrong:} Fraud is a source. \\[2pt]
		\textbf{C wrong:} IT failure is operational. \\[2pt]
		\textbf{D wrong:} Market loss is market risk.
	\end{frame}
	
	% ============ Q95 ============
	\begin{frame}[t]{Q95}
		\small
		\textbf{Question:} Which is the best description of "assurance engagement"?
		\begin{enumerate}[label=\Alph*.,itemsep=1pt,topsep=2pt]
			\item An objective assessment of evidence to provide an opinion or conclusion.
			\item A management advisory service.
			\item The marketing plan.
			\item The audit committee's schedule.
		\end{enumerate}
		\vspace{4pt}
		\textbf{Correct Answer: A}
	\end{frame}
	
	\begin{frame}[t]{Q95 -- Explanation}
		\small
		\textbf{Correct: A.} An assurance engagement objectively examines evidence to provide an independent opinion or conclusion on the subject matter. \\[4pt]
		\textbf{B wrong:} Advisory (consulting) is different. \\[2pt]
		\textbf{C wrong:} Marketing is unrelated. \\[2pt]
		\textbf{D wrong:} The committee's schedule is unrelated.
	\end{frame}
	
	% ============ Q96 ============
	\begin{frame}[t]{Q96}
		\small
		\textbf{Question:} Which is the best description of "consulting engagement"?
		\begin{enumerate}[label=\Alph*.,itemsep=1pt,topsep=2pt]
			\item An advisory service to add value without the auditor assuming management responsibilities.
			\item A financial statement audit.
			\item A compliance audit.
			\item An operational audit.
		\end{enumerate}
		\vspace{4pt}
		\textbf{Correct Answer: A}
	\end{frame}
	
	\begin{frame}[t]{Q96 -- Explanation}
		\small
		\textbf{Correct: A.} Consulting is advisory, adds value, and does not involve the auditor taking on management responsibilities. \\[4pt]
		\textbf{B wrong:} Financial statement audit is assurance. \\[2pt]
		\textbf{C wrong:} Compliance audit is assurance. \\[2pt]
		\textbf{D wrong:} Operational audit is assurance.
	\end{frame}
	
	% ============ Q97 ============
	\begin{frame}[t]{Q97}
		\small
		\textbf{Question:} Which is the best description of "follow-up"?
		\begin{enumerate}[label=\Alph*.,itemsep=1pt,topsep=2pt]
			\item Monitoring whether management actions were implemented.
			\item The audit report.
			\item The audit plan.
			\item The engagement letter.
		\end{enumerate}
		\vspace{4pt}
		\textbf{Correct Answer: A}
	\end{frame}
	
	\begin{frame}[t]{Q97 -- Explanation}
		\small
		\textbf{Correct: A.} Follow-up monitors whether management has implemented the agreed corrective actions, as required by Standard 2500. \\[4pt]
		\textbf{B wrong:} The report is a deliverable. \\[2pt]
		\textbf{C wrong:} The plan is a planning document. \\[2pt]
		\textbf{D wrong:} The engagement letter defines the agreement.
	\end{frame}
	
	% ============ Q98 ============
	\begin{frame}[t]{Q98}
		\small
		\textbf{Question:} Which is the best description of "quality assurance and improvement program" (QAIP)?
		\begin{enumerate}[label=\Alph*.,itemsep=1pt,topsep=2pt]
			\item A program to evaluate and improve the internal audit activity.
			\item The audit report.
			\item The audit plan.
			\item The engagement letter.
		\end{enumerate}
		\vspace{4pt}
		\textbf{Correct Answer: A}
	\end{frame}
	
	\begin{frame}[t]{Q98 -- Explanation}
		\small
		\textbf{Correct: A.} The QAIP covers all aspects of the internal audit activity, evaluating conformance with the Standards and the Code, and identifying improvement opportunities. \\[4pt]
		\textbf{B wrong:} The report is a deliverable. \\[2pt]
		\textbf{C wrong:} The plan is a planning document. \\[2pt]
		\textbf{D wrong:} The engagement letter defines the agreement.
	\end{frame}
	
	% ============ Q99 ============
	\begin{frame}[t]{Q99}
		\small
		\textbf{Question:} Which is the best description of "Code of Ethics"?
		\begin{enumerate}[label=\Alph*.,itemsep=1pt,topsep=2pt]
			\item Principles and rules of conduct for internal auditors.
			\item The audit report.
			\item The audit plan.
			\item The engagement letter.
		\end{enumerate}
		\vspace{4pt}
		\textbf{Correct Answer: A}
	\end{frame}
	
	\begin{frame}[t]{Q99 -- Explanation}
		\small
		\textbf{Correct: A.} The IIA Code of Ethics has two components: Principles (Integrity, Objectivity, Confidentiality, Competency) and Rules of Conduct that guide internal auditors' behavior. \\[4pt]
		\textbf{B wrong:} The report is a deliverable. \\[2pt]
		\textbf{C wrong:} The plan is a planning document. \\[2pt]
		\textbf{D wrong:} The engagement letter defines the agreement.
	\end{frame}
	
	% ============ Q100 ============
	\begin{frame}[t]{Q100}
		\small
		\textbf{Question:} Which is the best description of "internal audit charter"?
		\begin{enumerate}[label=\Alph*.,itemsep=1pt,topsep=2pt]
			\item A formal document defining the purpose, authority, and responsibility of internal audit.
			\item The audit report.
			\item The audit plan.
			\item The engagement letter.
		\end{enumerate}
		\vspace{4pt}
		\textbf{Correct Answer: A}
	\end{frame}
	
	\begin{frame}[t]{Q100 -- Explanation}
		\small
		\textbf{Correct: A.} The charter formally defines the internal audit activity's purpose, authority, and responsibility, and is approved by the board. \\[4pt]
		\textbf{B wrong:} The report is a deliverable. \\[2pt]
		\textbf{C wrong:} The plan is a planning document. \\[2pt]
		\textbf{D wrong:} The engagement letter is for a specific engagement.
	\end{frame}
	
	\begin{frame}{End}
		\centering
		\Large End of 100-Question Deck\\[8pt]
		\normalsize Review, repeat, and mark weak areas.\\[4pt]
		\small Add more questions by copying the Q/Explanation frame pairs.
	\end{frame}
	
\end{document}